\subsection{Intermediate currents of a shrinking sequence}\label{subsec:shrinking-defect-currents}

Let $Y$ be a connected compact complex manifold of dimension $n\geq1$, and let $h$ be a Hermitian form on $Y$. Consider smooth functions
\[
 v_j=s_ju_j,\qquad u_j\in\PSH(Y,h),\qquad
 \sup_Yu_j=0,\qquad 0<s_j\longrightarrow0.
\]
Suppose that their Monge--Amp\`ere masses equal a fixed finite number $b$ and that
\[
 (h+\ddc v_j)^n\rightharpoonup h^n+\sigma.
\]
The lower bound $h+\ddc v_j\geq(1-s_j)h$ makes $\sigma$ positive. Although $\ddc v_j\to0$ in the trace norm, its intermediate powers can retain a nonzero limit.

\begin{proposition}\label{prop:shrinking-intermediate-currents}
\index{Mass defect!intermediate currents}
After passage to a common subsequence, the positive forms
\[
 H_j=\ddc v_j+s_jh=s_j(h+\ddc u_j)
\]
satisfy $H_j^k\rightharpoonup K_k$ for $1\leq k\leq n$. The currents $K_k$ are positive and $\mathrm{d}$-closed, and $K_1=K_n=0$. For every smooth real $(n-k,n-k)$-form $\eta$ with $\mathrm{d}\eta=0$ or $\ddc\eta=0$, one has
\[
 \int_Y K_k\wedge\eta=0.
\]
Moreover,
\begin{equation}\label{eq:shrinking-defect-current-decomposition}
 \sigma=\sum_{k=2}^{n-1}\binom nk h^{n-k}\wedge K_k.
\end{equation}
If $\vartheta$ is a smooth closed semipositive $(1,1)$-form, then $\sigma$ is supported on $\{\vartheta^n=0\}$. In dimension three the same conclusion holds for a smooth semipositive form $\vartheta$ satisfying $\ddc\vartheta=0$.
\end{proposition}
\begin{proof}
For $s_j<1/2$, positivity in the expansion
\[
 (h+\ddc v_j)^n
 =\sum_{k=0}^n\binom nk(1-s_j)^{n-k}h^{n-k}\wedge H_j^k
\]
gives a uniform bound for each $\int_Yh^{n-k}\wedge H_j^k$. In local unitary coordinates, the coefficient matrix of $H_j^k$ on exterior powers is positive semidefinite. Its entries are bounded in absolute value by its diagonal trace, up to a dimensional constant. These trace bounds therefore bound every coefficient in total variation. Weak compactness of measures gives a common subsequence and positive limits $K_k$.

For $n\geq2$, choose a Gauduchon form $\gamma$. By \cref{prop:gauduchon-trace},
\[
 \int_Y H_j\wedge\gamma^{n-1}
 =s_j\int_Yh\wedge\gamma^{n-1}.
\]
Equivalence of the two smooth metrics implies $H_j\to0$ in coefficient total variation. For $n=1$ the same assertion follows from $\int_YH_j=s_j\int_Yh$. Thus $K_1=0$ in all dimensions. Since $\mathrm{d}H_j=s_j\mathrm{d}h$,
\[
 \mathrm{d}(H_j^k)=ks_j\mathrm{d}h\wedge H_j^{k-1}\longrightarrow0
\]
as currents; the preceding trace bounds control multiplication by the fixed smooth form $\mathrm{d}h$. Hence every $K_k$ is closed.

The smooth form $(\ddc v_j)^k$ is both $\mathrm{d}$-exact and $\ddc$-exact:
\[
 (\ddc v_j)^k
 =\mathrm{d}\bigl(\dc v_j\wedge(\ddc v_j)^{k-1}\bigr)
 =\ddc\bigl(v_j(\ddc v_j)^{k-1}\bigr).
\]
Its pairing with either kind of test form $\eta$ in the statement is zero by Stokes' formula. Expanding $\ddc v_j=H_j-s_jh$ gives
\[
 \int_Y\eta\wedge H_j^k
 =-\sum_{r=0}^{k-1}\binom kr(-s_j)^{k-r}
      \int_Y\eta\wedge h^{k-r}\wedge H_j^r
 =O(s_j).
\]
Each integral on the right is bounded by the trace bound for $H_j^r$. Taking $k=n$ and $\eta=1$ shows that the positive measures $H_j^n$ tend to zero in total variation. This proves $K_n=0$; the other pairings follow by weak convergence. Taking limits in the positive expansion proves \eqref{eq:shrinking-defect-current-decomposition}.

For closed semipositive $\vartheta$, apply the vanishing pairing to $\eta=\vartheta^{n-k}$. The measure $K_k\wedge\vartheta^{n-k}$ is positive and has zero mass. On every relatively compact open set where $\vartheta>0$, it controls the trace of $K_k$, so $K_k$ vanishes there. The decomposition gives the asserted support. In dimension three only $K_2$ occurs, and the same argument uses the $\ddc$-closed test form $\vartheta$ directly.
\end{proof}

In dimensions at most two the decomposition forces $\sigma=0$. A threefold carrying a pluriclosed Hermitian metric also has no positive excess defect of this form. In higher dimensions, the vanishing pairings impose necessary conditions on each intermediate current. They do not by themselves construct a smooth sequence realizing prescribed currents.
